% Options for packages loaded elsewhere
\PassOptionsToPackage{unicode}{hyperref}
\PassOptionsToPackage{hyphens}{url}
\documentclass[
]{article}
\usepackage{xcolor}
\usepackage{amsmath,amssymb}
\setcounter{secnumdepth}{-\maxdimen} % remove section numbering
\usepackage{iftex}
\ifPDFTeX
  \usepackage[T1]{fontenc}
  \usepackage[utf8]{inputenc}
  \usepackage{textcomp} % provide euro and other symbols
\else % if luatex or xetex
  \usepackage{unicode-math} % this also loads fontspec
  \defaultfontfeatures{Scale=MatchLowercase}
  \defaultfontfeatures[\rmfamily]{Ligatures=TeX,Scale=1}
\fi
\usepackage{lmodern}
\ifPDFTeX\else
  % xetex/luatex font selection
\fi
% Use upquote if available, for straight quotes in verbatim environments
\IfFileExists{upquote.sty}{\usepackage{upquote}}{}
\IfFileExists{microtype.sty}{% use microtype if available
  \usepackage[]{microtype}
  \UseMicrotypeSet[protrusion]{basicmath} % disable protrusion for tt fonts
}{}
\makeatletter
\@ifundefined{KOMAClassName}{% if non-KOMA class
  \IfFileExists{parskip.sty}{%
    \usepackage{parskip}
  }{% else
    \setlength{\parindent}{0pt}
    \setlength{\parskip}{6pt plus 2pt minus 1pt}}
}{% if KOMA class
  \KOMAoptions{parskip=half}}
\makeatother
\usepackage{longtable,booktabs,array}
\newcounter{none} % for unnumbered tables
\usepackage{calc} % for calculating minipage widths
% Correct order of tables after \paragraph or \subparagraph
\usepackage{etoolbox}
\makeatletter
\patchcmd\longtable{\par}{\if@noskipsec\mbox{}\fi\par}{}{}
\makeatother
% Allow footnotes in longtable head/foot
\IfFileExists{footnotehyper.sty}{\usepackage{footnotehyper}}{\usepackage{footnote}}
\makesavenoteenv{longtable}
\setlength{\emergencystretch}{3em} % prevent overfull lines
\providecommand{\tightlist}{%
  \setlength{\itemsep}{0pt}\setlength{\parskip}{0pt}}
\usepackage[T2A]{fontenc}
\usepackage[russian]{babel}
\usepackage{paratype}
\clubpenalty=10000
\widowpenalty=1000
\DeclareUnicodeCharacter{0301}{\hspace{-1ex}\'{ }}
\usepackage{bookmark}
\IfFileExists{xurl.sty}{\usepackage{xurl}}{} % add URL line breaks if available
\urlstyle{same}
\hypersetup{
  pdftitle={Указатель заглавий произведений художественной литературы{,} 1801--1975},
  pdfauthor={Кирилл Маслинский, Евгения Лекаревич},
  hidelinks,
  pdfcreator={LaTeX via pandoc}}

\title{Указатель заглавий произведений художественной литературы,
1801--1975}
\author{Кирилл Маслинский, Евгения Лекаревич}
\date{}

\begin{document}
\maketitle
\begin{abstract}
В датасете в машиночитаемом виде представлены тома «Указателя заглавий
произведений художественной литературы, 1801--1975». Печатный указатель
предназначен для определения имени автора по заглавию художественного
произведения. В нем представлены художественные произведения, изданные в
виде отдельных книг, на русском языке и в переводах с иностранных языков
и языков народов СССР с~1801 по~1975 год. Включены издания «вольной
печати» и зарубежных издательств, сотрудничавших с РСФСР в 1920-е гг.
Для удобства пользователей записи из указателей были разделены на
отдельные поля (автор, заглавие, альтернативное заглавие) с помощью
автоматического анализатора. Данные представлены в табличной форме.
\end{abstract}

\begin{quote}
Формат цитирования датасета: \emph{Маслинский, Кирилл; Лекаревич,
Евгения} Указатель заглавий произведений художественной литературы,
1801--1975 // Репозиторий открытых данных по русской литературе и
фольклору. 2026. V1. https://doi.org/10.31860/openlit-2022.12-B012
\end{quote}

\subsection{Общее описание
данных}\label{ux43eux431ux449ux435ux435-ux43eux43fux438ux441ux430ux43dux438ux435-ux434ux430ux43dux43dux44bux445}

Машиночитаемая библиографическая база данных отражает заглавия всех
видов, родов и жанров художественной литературы. Базируется на
семитомном «Указателе заглавий произведений художественной литературы».
В указателе представлены книги, написанные на русском языке, и издания
переводов с иностранных языков и языков народов СССР. В «Указатель»
включено около 130000 заглавий художественных произведений,
опубликованных в виде отдельных книг на территории дореволюционной
России и СССР (в том числе издания «вольной русской печати», а также
зарубежных издательств, сотрудничавших с РСФСР в 1920-е гг.). Для
удобства пользователей записи из указателей были разделены на отдельные
поля (автор, заглавие, альтернативное заглавие) с помощью
автоматического анализатора. Данные представлены в табличной форме.

«Алфавитный указатель заглавий художественных произведений» выделяет из
потока книг, вышедших в дореволюционной России и СССР, произведения
\emph{художественной} литературы, включая иностранные сочинения,
переводившиеся на русский язык в течение 175 лет. Работа над этим
изданием велась коллективами специалистов из Государственной библиотеки
СССР им. В. И. Ленина, Государственной публичной библиотеки имени
М.~Е.~Салтыкова-Щедрина и Библиотеки Академии наук СССР, начиная
с~7~тома --- РГБ, РНБ, БАН, основой для составления «Указателя»
послужили генеральные алфавитные каталоги этих библиотек. С критериями
отбора произведений для указателя можно ознакомиться в предисловии от
составителей к первому тому (с. 7--10). В настоящем датасете
представлены шесть томов указателя, охватывающие период 1801--1975 гг.
Седьмой том, содержащий указатель всех авторов, упомянутых в первых
шести томах, в настоящий датасет не включен.

Указатель заглавий произведений художественной литературы, 1801--1975 :
{[}в 8 т.{]} / Гос. б-ка СССР им. В.И. Ленина, Гос. публ. б-ка
им.~М.~Е.~Салтыкова-Щедрина, Б-ка Акад. наук СССР ; сост.:
Б.~М.~Вишневецкая, А.~А.~Колганова \ldots{} Е.~К.~Соколинский
{[}и~др.{]}. -- Москва, 1985--1995.

\begin{itemize}
\tightlist
\item
  Т. 1 : А-Г. -- 1985. -- 488 с.
\item
  Т. 2 : Д-З. -- 1986. -- 349 с.
\item
  Т. 3 : И-Л. -- 1989. -- 326 с.
\item
  Т. 4 : М-О. -- 1990. -- 431 с.
\item
  Т. 5 : П-С. -- 1990. -- 477 с.
\item
  Т. 6 : С(Сел.) - Я. -- 1991. -- 544 с.
\item
  Т. 7 : {[}Указатели{]} / Рос. гос. б-ка, Рос. нац. б-ка, Б-ка Акад.
  наук Рос. Федерации. -- 1995. -- 536 с.
\end{itemize}

Данные были получены в результате автоматизированного парсинга
отсканированных указателей. Работа велась итеративно: по результатам
парсинга проводилась проверка и правка ошибок в исходных текстовых
файлах указателей, корректировка алгоритмов парсера и запуск парсера по
откорректированным данным. Все исправления вносятся в исходные текстовые
файлы или в код парсера, таблицы данных целиком порождаются парсером.
Массовые исправления (переносы, ошибочно распознанные символы, разрывы в
именах авторов) выполнены сценариями из каталога \texttt{scripts/} с
использованием указателя имен авторов (т.~7) для проверки написания
фамилий; ручные исправления, сверенные со сканами, сгруппированы в
истории изменений репозитория по видам ошибок.

Текущая версия всех исходных файлов и код, на основании которых
составлена база, находится в гит-репозитории по адресу:
https://github.com/maslinych/fiction-titles. Авторы датасета будут
признательны за сообщения о замеченных ошибках и неточностях. Для связи
можно воспользоваться контактами, указанными в метаданных датасета. По
мере исправления ошибок и дополнения данных будут публиковаться новые
версии датасета.

\subsection{Состав
датасета}\label{ux441ux43eux441ux442ux430ux432-ux434ux430ux442ux430ux441ux435ux442ux430}

\begin{itemize}
\tightlist
\item
  \texttt{vol\_1.csv\ ...\ vol\_6.csv} --- списки заглавий произведений
  (по одному файлу на том указателя).
\item
  \texttt{README.md} --- этот файл.
\end{itemize}

\subsection{Сводка}\label{ux441ux432ux43eux434ux43aux430}

{\def\LTcaptype{none} % do not increment counter
\begin{longtable}[]{@{}
  >{\raggedright\arraybackslash}p{(\linewidth - 14\tabcolsep) * \real{0.1622}}
  >{\raggedright\arraybackslash}p{(\linewidth - 14\tabcolsep) * \real{0.2027}}
  >{\raggedleft\arraybackslash}p{(\linewidth - 14\tabcolsep) * \real{0.1351}}
  >{\raggedleft\arraybackslash}p{(\linewidth - 14\tabcolsep) * \real{0.1216}}
  >{\raggedleft\arraybackslash}p{(\linewidth - 14\tabcolsep) * \real{0.0541}}
  >{\raggedleft\arraybackslash}p{(\linewidth - 14\tabcolsep) * \real{0.0541}}
  >{\raggedleft\arraybackslash}p{(\linewidth - 14\tabcolsep) * \real{0.1757}}
  >{\raggedleft\arraybackslash}p{(\linewidth - 14\tabcolsep) * \real{0.0946}}@{}}
\toprule\noalign{}
\begin{minipage}[b]{\linewidth}\raggedright
Файл
\end{minipage} & \begin{minipage}[b]{\linewidth}\raggedright
Буквы
\end{minipage} & \begin{minipage}[b]{\linewidth}\raggedleft
Записей
\end{minipage} & \begin{minipage}[b]{\linewidth}\raggedleft
alt\_title
\end{minipage} & \begin{minipage}[b]{\linewidth}\raggedleft
in
\end{minipage} & \begin{minipage}[b]{\linewidth}\raggedleft
note
\end{minipage} & \begin{minipage}[b]{\linewidth}\raggedleft
Непустой tail
\end{minipage} & \begin{minipage}[b]{\linewidth}\raggedleft
MISSING
\end{minipage} \\
\midrule\noalign{}
\endhead
\bottomrule\noalign{}
\endlastfoot
vol\_1.csv & А--Г & 19041 & 591 & 334 & 9 & 221 (1.2\%) & 0 \\
vol\_2.csv & Д--З & 13249 & 389 & 184 & 11 & 76 (0.6\%) & 0 \\
vol\_3.csv & И--Л & 12475 & 361 & 296 & 10 & 64 (0.5\%) & 0 \\
vol\_4.csv & М--О & 17013 & 458 & 420 & 18 & 99 (0.6\%) & 0 \\
vol\_5.csv & П--С (Сек.) & 19146 & 591 & 397 & 12 & 96 (0.5\%) & 0 \\
vol\_6.csv & С (Сел.)--Я, A--Z & 22888 & 627 & 490 & 11 & 256 (1.1\%) &
1 \\
\textbf{Всего} & & \textbf{103812} & 3017 & 2121 & 71 & 812 (0.8\%) &
1 \\
\end{longtable}
}

Для каждого файла указаны число записей и число записей с непустыми
полями \texttt{alt\_title}, \texttt{in}, \texttt{note} и \texttt{tail},
а также число пропущенных номеров (\texttt{MISSING}). Таблица
порождается командой \texttt{make\ stats} (сценарий
\texttt{scripts/stats.py}).

\subsection{Формат
данных}\label{ux444ux43eux440ux43cux430ux442-ux434ux430ux43dux43dux44bux445}

\subsubsection{vol\_N.csv}\label{vol_n.csv}

Колонки в таблице данных:

\begin{itemize}
\tightlist
\item
  \texttt{start} --- номер строки в исходном txt-файле
  (\texttt{txt/vol\_N.txt} в репозитории), на которой начинается запись.
\item
  \texttt{end} --- номер строки в исходном txt-файле, на которой
  оканчивается запись (включая следующие за записью пустые строки).
\item
  \texttt{num} --- порядковый номер записи, как в источнике. Некоторые
  номера повторяются и могут быть снабжены буквенными индексами (1б),
  другие пропущены, в этих случаях вместо одного номера записи указан
  диапазон (101---102). Если номер пропущен в самом источнике, в таблице
  есть строка с этим номером, где \texttt{title} = \texttt{NOPARSE} и
  \texttt{tail} = \texttt{MISSING} (единственный случай --- т.~6,
  №~19080).
\item
  \texttt{title} --- заглавие произведения, включая подзаголовок
  (отделен двоеточием).
\item
  \texttt{author} --- автор(ы) издания. Формат указания автора --- как в
  источнике: фамилии авторов выделены заглавными буквами, раскрытие
  псевдонима дано в скобках. Если автор не установлен, указано «Имя авт.
  не установлено.»; у нескольких записей, где в источнике нет сведений
  об авторе, поле пустое.
\item
  \texttt{alt\_title} --- альтернативные названия, под которыми
  издавалось произведение (пометы «Изд. также под загл.:», «На тит. л.
  загл.:», «Загл. обл.:»).
\item
  \texttt{in} --- заглавие издания, в составе которого напечатано
  произведение (помета «В кн.:»).
\item
  \texttt{note} --- примечания об атрибуции: «Ошибочно
  приписывалось\ldots», «Предполагаемое имя авт.: \ldots».
\item
  \texttt{tail} --- фрагмент записи, который не удалось автоматически
  распределить по остальным колонкам. Здесь могут отражаться ошибки
  распознавания либо нестандартные пометы, использованные в Указателе
  (например, «1-я пьеса изд. также под загл.: \ldots»).
\end{itemize}

\end{document}
